\subsection{Two-party $L_1$ Sampling}\label{sec:l1}
In this section, we describe a secure two-party $L_1$ sampling protocol.
Given two $n$-dimensional vectors $\vec{w}_1 = (w_{1,1}, \ldots,
w_{1,n})$ and $\vec{w}_2 = (w_{2,1}, \ldots, w_{2,n})$ as the private
inputs from parties $P_1$ and $P_2$ respectively, the protocol samples from the $L_1$ distribution according to
$\vec{w}_1 + \vec{w}_2$.  

\paragraph{Notation: $L_p$ norm.}
Let $\vec{w} = (w_1, \ldots, w_n) \in \RR^n$ be a non-zero vector. \emph{The $L_p$
norm $\Vert \w \Vert_p$ of $\vec{w}$} is defined as
$\Vert \w \Vert_p := \left ( \sum_j |w_j|^p \right )^{1/p} .$
When there is no subscript, it means $L_2$ norm; that is, $\Vert \w \Vert := \Vert \w \Vert_2$  

\paragraph{Assumptions.}
Throughout the paper, we assume that the values $w_{b,i}$ are
represented by fixed-point precision numbers, and consider the cost of
communicating such a number to be independent of $n$.
% \anote{This is the
% only place it felt appropriate to say something about precision.  Do you
% like what I wrote?}
%
We assume all weights in vectors $\vec{w}_1$ and $\vec{w}_2$ are
non-negative.

\paragraph{Ideal functionality.}
We first define an ideal functionality for the two-party $L_1$ sampling.
Slightly abusing the notation, let $L_1(\vec{w}_1, \vec{w}_2)$ be a
two-input sampling procedure based on the $L_1$ distribution of $\vec{w}_1 +
\vec{w}_2$: 
%
$$\Pr[ L_1(\vec{w}_1, \vec{w}_2) \mbox{ samples } i ] = \frac{w_{1,i} +
w_{2,i}}{ \|\vec{w}_1 + \vec{w}_2\|_1 }. $$


We give a more formal description of the functionality $\F_{L_1}$ in the
figure below.  In Section \ref{sec:L1prot}, we present a protocol that
securely realizes this functionality. 

\bbox
{
\begin{center}
 $\F_{L_1}$:  Ideal functionality for two-party $L_1$ sampling
\end{center}

The functionality has the following parameter:
\BI
\item $n \in \NN$. The dimension of the input weight vectors $\vec{w}_1$ and $\vec{w}_2$.
\EI

The functionality proceeds as follows:
\BEN
\item Receive inputs $\vec w_1$ and $\vec w_2$ from $P_1$ and $P_2$
respectively. 
\item Sample $i \in [n]$ with probability $\frac{w_{1,i} +
w_{2,i}}{ \|\vec{w}_1 + \vec{w}_2\|_1 }$
\item Send $i$ to $P_1$ and $P_2$. 

\EEN 
}

\subsubsection{A toy protocol towards securely realizing $\F_{L_1}$}
We describe our first attempt, which is  insecure, but provides good
intuition on how we construct a secure protocol.  In fact, the attack on
this broken protocol, as well as the fix presented in the next
sub-section, remain relevant when we move to product sampling and $L_2$
sampling as well.  Since we assume that all the weights are
non-negative, we observe that letting $p = \frac{\|\vec{w}_1\|_1}{
\|\vec{w}_1\|_1 + \|\vec{w}_2\|_1}$, the above measure can be re-written
as follows:

\begin{equation}
\label{eqn:l1}
\Pr[ L_1(\vec{w}_1, \vec{w}_2) \mbox{ samples } i ] = 
  \frac{w_{1,i}}{\|\vec{w}_1\|_1} \cdot p  + \frac{w_{2,i}}{\|\vec{w}_2\|_1} \cdot (1-p).
\end{equation}


Equation (\ref{eqn:l1}) leads us to the following natural approach.

\BEN
\item Party $P_1$ samples $i_1$ from the $L_1$ distribution according to
$\vec{w}_1$, such that  $\Pr[P_1 \mbox{ samples } i_1] =
\frac{w_{1,i_1}}{\|\vec{w}_1\|_1}$.

\item Party $P_2$ samples $i_2$ from the $L_1$ distribution according to $\vec{w}_2$,
such that $\Pr[P_2 \mbox{ samples } i_2] =
\frac{w_{2,{i_2}}}{\|\vec{w}_2\|_1}$.

\item Then, $P_1$ and $P_2$ execute a secure protocol for the following
procedure:
  \BEN
  \item Execute a coin toss protocol with bias $p$.  Let $b$ be the output of the coin-flip. 
  \item If $b=0$ (resp., $b=1$), output $i_1$ (resp., $i_2$).
  \EEN
\EEN

The output of the protocol will achieve correct sampling.  

\paragraph{Insecurity of the protocol.} However, this protocol has a
subtle security issue.
%
For example, let $i$ be the eventual output index of the protocol. Then,
we have the following:
%
\BI
\item If the coin flip $b$ is 0, which happens with probability $p$, it
holds that $i$ is always the same as $i_1$. 

\item On the other hand, if the coin flip $b$ is 1, then $i$ will be
the same as $i_1$ if and only if $i_2 = i_1$, which happens with
probability $\frac{w_{2, i_1}}{ \| \vec w_2 \|_1 }$.
\EI
%
This implies that we have 
$$\Pr[ i = i_1 | i_1 ] = p + (1-p) \cdot \frac{w_{2, i_1}}{ \| \vec w_2 \|_1 }$$

Now consider a distinguisher that corrupts $P_1$, chooses inputs $\w_1$
and $\w_2$, and checks the above conditional probability, which is
possible {\em since the distinguisher can also see $i_1$ through the corrupted
$P_1$}.
%
To prove security, we should be able to construct a simulator for $P_1$
that fools this distinguisher. However, a simulator for $P_1$ doesn't know
$\w_2$, which causes the above conditional probability to be
unsimulatable. 
%
%Interpreting this, if the coin flip $b$ is 0, which
%happens with probability $p$, it holds that $i$ is always the same as
%$i_1$. On the other hand, if the coin flip $b$ is 1, then $i$ will be
%the same as $i_1$ if and only if $i_2 = i_1$, which happens with
%probability $\frac{w_{2, i_1}}{ \| \vec w_2 \|_1 }$.

In a sense, by having $P_1$ choose $i_1$, the protocol allows $P_1$ to
measure the conditional probability $\Pr[ i = i_1 | i_1 ]$, which
depends on the value $w_{2, i_1}$ thereby leaking information about
$P_2$'s input to $P_1$.   


\subsubsection{Secure $L_1$ sampling protocol}  \label{sec:L1prot}
\paragraph{Oblivious sampling.}  
We address the insecurity of the toy protocol by having the parties
sample \emph{obliviously} from $\vec w_1$, $\vec w_2$.  This way, each
party would not know whether the final output index matches the sample
taken from its own vector, or the sample taken from the other party's
vector. Specifically, we will construct our protocol under the framework
described below: 

\BEN
\item The parties obliviously sample $i_1$ according to $L_1$
distribution of $\vec w_1$. The output index $i_1$ is secret shared
between the two parties. Let $\langle i_1 \rangle$ denote the secret share of $i_1$.
Likewise, they obliviously sample $\langle i_2 \rangle$ from $L_1$ distribution of
$\vec w_2$.

\item Execute a secure two-party protocol to compute the following:
  \BEN
  \item Flip a coin $b$ with bias $p$. 
  \item If $b=0$, output the decryption of $i_1$; otherwise output
  the decryption of $i_2$. 
  \EEN
\EEN


\paragraph{Ideal functionalities.}
Formally, we define an ideal functionality $\F_{\oslone}$ as follows:

\bbox
{
\begin{center}
 $\F_{\oslone}$:  Ideal functionality for oblivious $L_1$ sampling.
\end{center}

The functionality considers two participants, the sender and the
receiver. The functionality is parameterized with a number $n$. 

\medskip

\textbf{Inputs:} The sender has an $n$-dimensional weight vector $\w$.
The receiver has no input. 

\medskip

The functionality proceeds as follows:
\BEN
\item Receive $\vec w$ from the sender.
\item Sample $i \in [n]$ with probability $\frac{w_i}{ \|\vec{w}\|_1 }$
\item Choose a random pad $\pi \in \zo^\ell$, where $\ell = \lceil \log_2 n \rceil$. 
\item Send $\pi$ to the sender and $i \xor \pi$ to the receiver. 

\EEN 
}

\medskip
We also give an ideal functionality $\F_{\bcoin}$ for the biased coin tossing. 

\bbox
{
\begin{center}
 $\F_{\bcoin}$:  Ideal functionality for biased coin tossing.
\end{center}

The functionality considers two participants $P_1$ and $P_2$ and proceeds as
follows:

\BEN
\item Receive a number $s_1$ as input from $P_1$ and $s_2$ from $P_2$.  
\item Flip a coin $b$ with bias $p = \frac{s_1}{s_1 + s_2}$. 
\item Choose a random bit $r \in \zo$.
\item Send $r$ to $P_1$ and $r \xor b$ to the receiver. 
\EEN 
}

\medskip

\paragraph{$L_1$ sampling protocol.}
Based on the above functionalities, we describe a protocol securely realizing
$\F_{L_1}$ in the $(\F_\oslone, \F_\bcoin)$-hybrid. 

\floatname{algorithm}{Protocol}
\begin{algorithm}[h]
\textbf{Inputs:} Party $P_b$ has input $\vec{w}_b$. 

\BEN
\item Execute $\F_\oslone$ with $P_1$ as a sender with input $\vec{w}_1$
and $P_2$ as a receiver. Let $\langle i_1 \rangle$ be the secret share of the output
index. 

\item Execute $\F_\oslone$ with $P_2$ as a sender with input $\vec{w}_2$
and $P_1$ as a receiver. Let $\langle i_2 \rangle$ be the secret share of the output
index.

\item Execute $\F_\bcoin$ where $P_1$ has input $\|\vec{w}_1\|_1$ and
$P_2$ has input $\|\vec{w}_2\|_1$. Let $\langle b \rangle$ be the secret share of the
output bit.

\item Execute $\F_{2PC}$ for the following circuit: 
      \BEN
      \item Input: $\langle i_1 \rangle, \langle i_2 \rangle, \langle b \rangle$. 
      \item Output: $i_1 \cdot (1-b) + i_2 \cdot b$.
      \EEN
\EEN

\caption{Two-party $L_1$ sampling in the $(\F_\oslone, \F_\bcoin)$-hybrid.}
\label{fig:l1}
\end{algorithm}

\begin{theorem} \label{th:l1}
Protocol~\ref{fig:l1} securely realizes $\F_{L_1}$ with semi-honest
security in the $(\F_\oslone, \F_\bcoin)$-hybrid.
\end{theorem} 
The proof is found in Section~\ref{apx:pfTheorem1}.


\paragraph{Securely realizing $\F_\oslone$ with threshold FHE.}
%In the 1-out-of-$m$ OT-based sampling protocol, the weights have to fit
%into $m$ slots of the OT sender input. This level of precision can often
%be acceptable depending on the circumstances, but there will be
%situations where a higher level of precision is necessary.  Here, we
%describe a sampling protocol based on a threshold FHE scheme that
%provides a higher level of precision.
%
The main idea of the protocol is having the parties securely sample a
random number $r$ from $[s]$, where $s := \| \vec{w} \|_1$. Our construction is found in
Protocol~\ref{fig:os-fhe}.

\iffalse
, which can be performed by the following
circuit.

\bbox
{
\BEN
\item[] \hspace{-2em} $\rhd$ Goal: Sample a random number from $[s]$
\item[] \hspace{-2em} $\rsample( s, \{r_{1,j}\}_{j=1}^\secparam, \pi_1, \{r_{2,j}\}_{j=1}^\secparam, \pi_2)$   ~~~$\rhd$ $r_{b,j}, \pi_b$ will
be all random bits. 

\item The sender's input is $(s, \{r_{1,j}\}_{j=1}^\secparam, \pi_1)$
and the receiver's input is $(\{r_{2,j}\}_{j=1}^\secparam, \pi_2)$.\\ We
assume that  all objects are represented with $\secparam$ bits.

\item Compute the tight bitmask $\beta = 0^{\ell-h}1^h$ such that $s \& \beta = s$ and $s
| \beta = \beta$ where $\&$ (resp., $|$) denotes bitwise AND (resp.,
bitwise OR) operation. Note that we have $h$ such that $2^{h-1} \le s <
2^h$. This computation can be done by checking all possible candidates
of $h$ one by one.

\item Let $\pi = \pi_1 \xor \pi_2$. For $j=1, \ldots, \secparam$, compute $r_j := (r_{1,j} \xor r_{2,j}) \& \beta$
\item Let $j^*$ be the first index such that $r_{j^*} \le s$. 
\item Output $\pi$ to the sender and output $r_{j^*} \xor \pi$ to the receiver.
\EEN
}

\medskip
Note that $\Pr[(r_j \& \beta) > s] < 1/2$ and therefore with
probability $(1 - 2^{-\secparam})$ there is a good $j^*$.
Based on this building block, we can achieve a oblivious sampling
protocol based on threshold FHE. See Protocol \ref{fig:os-fhe}.
\fi

\floatname{algorithm}{Protocol}
\begin{algorithm}[htbp]

\textbf{Inputs:} The sender has input $\vec{w} = (w_1, \ldots, w_n)$. 
\BEN
    \item The sender computes $s := \| \vec{w} \|_1$. 
    \item The sender and the receiver execute $\F_{2PC}$ to uniformly 
    sample $r$ from the range $[0, s)$. This is possible, since $s$ has
    a fixed point representation. Let $r_1$ and $r_2$ be the secret
    share of $r$ given to $P_1$ and $P_2$ respectively.  

    \item The sender and the receiver set up a threshold FHE scheme. The plaintext
    space of the FHE is $GF(2)$, which allows homomorphic bitwise-xor and
    bitwise-AND operations. Let $\lbra m \lket$ denote an FHE encryption
    of plaintext $m$ which can be a bit or bits depending on the
    context.

    \item The receiver sends $\lbra r_2 \lket$ so that the sender can
    compute $\lbra r \lket := \lbra r_1 \lket \xor \lbra r_2 \lket$. 

    \item The sender homomorphically evaluates the following circuit: 
        \BEN
        \item Let $cnt_0 = 0$. For $j=1,...,n$, let $cnt_j = cnt_{j} + w_{j}$. 
        \item Output $i \in [1,n]$ such that $r \in [cnt_{i-1}, cnt_{i}]$.
        \EEN 
        Let $\lbra i \lket$ be the output encryption from the above
        homomorphic evaluation. 
        
    \item The sender chooses a random pad $\pi$, and then it sends $\lbra c \lket =
    \lbra i \lket \xor \lbra \pi \lket$ to the receiver. 
    \item The two parties perform threshold decryption so that $c$ is
    decrypted to the receiver. 
    \item The sender outputs $\pi$ and the receiver outputs the
    decryption of $c$. 
\EEN
\caption{Oblivious sampling from threshold FHE}
\label{fig:os-fhe}
\end{algorithm}

\begin{theorem} \label{thm:os-fhe}
Assuming the existence of threshold FHE with
$\mathsf{IND}\mbox{-}\mathsf{CPA}$ security,
Protocol~\ref{fig:os-fhe} securely realizes $\F_{\oslone}$ in the
semi-honest security model. 
\end{theorem} 
The proof is found in Section~\ref{apx:pfTheorem2}.

We note that we give another construction that relies on sub-linear
1-out-of-$m$ oblivious transfer (OT), but requires computation that is
exponential in the bit precision in Section~\ref{apx:otl1}.  

\paragraph{Securely realizing $\F_\bcoin$.} The secure construction for
$\F_\bcoin$ is straightforward and can be found in Section~\ref{apx:bcoin}. 

